% ============================================================
%  THE ORTYX PROTOCOL
%  Part IV — Reconstruction
%  Chapter 18: Memory, Interference, and Field Dynamics
% ============================================================

\chapter{Memory, Interference, and Field Dynamics}
\label{ch:memory-field}

\chapepigraph{The question is not whether memory is a field.
The question is whether treating it as a field generates
predictions that treating it as storage does not. If it
does, the field is not a metaphor. It is a hypothesis.}{Flyxion}

\noindent
The \MEMeight{} architecture's translation into
RSVP-compatible terms is the most technically demanding
of the reconstruction tasks, because the \MEMeight{}
formalism is the most elaborate of the Phoenix corpus's
formal apparatus. The reconstruction works at two levels:
the mathematical level, where the wave superposition
formalism connects to RSVP field dynamics, and the
conceptual level, where the dissolution of the
storage--process distinction connects to the RSVP
principle of constraint-governed trajectory evolution.

\section{Wave Superposition as Field Dynamics}
\label{sec:wave-field}

The \MEMeight{} global memory field:
\[
  \MemField(x,t) = \sum_{i=1}^{N} M_i(x,t)
  = \sum_{i=1}^{N} A_i(x,t)\,e^{i(\omega_i t + \phi_i(x,t))}
\]
is a superposition of complex-valued wave functions
over a representational space $x$ and time $t$. In
RSVP-compatible terms, this can be read as a
parameterization of the RSVP scalar field $\phi(x,t)$
in the specific case where the field's dynamics are
governed by wave superposition rather than by the
general RSVP field equations.

The reading is metaphorical at the current state of
the reconstruction: it does not follow from the RSVP
field equations that the field should be a superposition
of complex exponentials with the specific amplitude
decay and emotional modulation terms of the \MEMeight{}
formalism. But the reading is productive: it suggests
that the \MEMeight{} formalism might be derivable as
a special case of the RSVP dynamics under specific
boundary conditions, which is a testable theoretical
claim.

What is not metaphorical is the interference energy:
\[
  E_{\mathrm{int}} = |\MemField|^2
  = \sum_i |M_i|^2 + \sum_{i \neq j} M_i\overline{M_j}.
\]
The cross-term $\sum_{i \neq j} M_i\overline{M_j}$
encodes $O(N^2)$ pairwise relationships at $O(N)$
storage cost. This is a mathematical fact about wave
superposition, independent of the interpretation of
$x$, $\omega_i$, or $\phi_i$. It is at \LevelI{} and
requires no RSVP translation to be established. The
RSVP translation adds the interpretation --- this is
a form of constraint-governed associative memory ---
without changing the mathematics.

\section{Memory as Constraint-Governed Trajectory}
\label{sec:memory-trajectory}

The more significant reconstruction is conceptual.
The \MEMeight{} architecture proposes that memory is
not a stored object but an ongoing process: the global
field $\MemField(x,t)$ is simultaneously the stored
information and the ongoing dynamics. This dissolution
of the storage--process distinction is, in RSVP terms,
the claim that memory is a constraint-governed trajectory
rather than a state.

In RSVP dynamics, the trajectory of the scalar field
$\phi(x,t)$ is governed by constraint equations that
determine how the field can evolve. The field's current
value at each point constrains its future values;
the constraints are the memory. Retrieving a memory,
in this frame, is not accessing a stored object but
allowing the constrained field dynamics to evolve
toward the attractor corresponding to the queried
pattern.

This RSVP framing makes explicit what the \MEMeight{}
architecture implies: the constraint structure of the
field --- the pattern of constructive and destructive
interference that determines which resonances are
amplified --- is the memory. Not the field values
themselves, which change continuously, but the pattern
of constraints that shapes how the field evolves.

\begin{salvage}[Reconstruction: Retrieval as Constraint Relaxation]
In the RSVP frame, memory retrieval is the application
of a query that partially specifies the field's
admissible trajectory space. The query $Q$ constrains
which field trajectories are admissible; the system
then evolves toward the constraint-satisfying trajectory
that is closest to the current field state. The
``resonance'' in the \MEMeight{} retrieval integral
$R_i(Q) = |\int Q(x)\overline{M_i(x,t)}\,dx|$ is
the degree to which memory $M_i$ is consistent with
the query's constraints. High resonance indicates
high constraint-compatibility; the memory with highest
resonance is the one whose admissible trajectory space
most overlaps with the query's constraints.
\end{salvage}

\section{The Storage-Process Dissolution in RSVP}
\label{sec:storage-process-rsvp}

The dissolution of the storage--process distinction
--- one of the philosophically most interesting features
of the \MEMeight{} architecture --- is a natural
consequence of the RSVP field-theoretic framework.
In RSVP, the scalar field $\phi(x,t)$ evolves according
to its constraint equations; there is no separate
``storage medium'' that holds values independently
of the dynamics. The field's current state is
simultaneously what the system knows (the stored
information, represented as constraint patterns) and
what the system is doing (the ongoing evolution under
those constraints).

This is not a unique feature of RSVP or \MEMeight{};
it is characteristic of any field-theoretic representation.
The electromagnetic field, the gravitational field,
and quantum fields all exhibit the same property:
there is no separate ``stored value'' independent of
the field dynamics. The \MEMeight{} contribution is
to propose that a similar field-theoretic structure
should govern the representation of memory in cognitive
systems, and the RSVP translation situates this
proposal within a broader theoretical context that
has developed the field-theoretic approach in other
domains.

\section{Emotional Modulation as Field Coupling}
\label{sec:emotional-modulation}

The emotional modulation of amplitude envelopes in
\MEMeight{}:
\[
  A_i(x,t) = A_i^0\exp(-\lambda_i t)(1 + \eta\,e(t))
\]
becomes, in the RSVP translation, a coupling between
the memory field and an external field representing
the system's affective state. The emotional coupling
parameter $\eta$ determines how strongly the affective
field modulates the memory field's dynamics; the
decay constant $\lambda_i$ represents the rate at
which memory $M_i$'s contribution to the field
dissipates in the absence of reinforcement.

This translation does not resolve the operationalization
gap identified in Chapter~12: $e(t)$ remains without
a measurement procedure. What the translation does
is situate the operationalization gap within a more
general framework: the problem of specifying $e(t)$
is the problem of operationalizing the affective field
coupling in a RSVP-compatible system, which connects
to existing work in computational affective computing
rather than being unique to the Phoenix corpus.

\section{The Biological Correspondence Hypothesis}
\label{sec:biological-correspondence}

Chapter~4 identified the biological correspondence
hypothesis --- the claim that the Marine--\MEMeight{}--\HCFone{}
architecture corresponds to known properties of
biological auditory and cognitive systems --- as a
testable prediction that the Phoenix corpus had not
yet systematically evaluated. The RSVP translation
does not provide the missing evaluation, but it
provides a more precise statement of what the hypothesis
predicts.

In RSVP-compatible terms, the biological correspondence
hypothesis claims:

(1)~The auditory system implements something functionally
equivalent to the Marine jitter metric in its
pre-attentive processing of sound: it preferentially
represents signals in low-$\AccEnt$ regions of the
acoustic trajectory space.

(2)~The hippocampal memory system implements something
functionally equivalent to \MEMeight{} wave interference
in its encoding and retrieval of episodic memories:
memories are represented as interference patterns in
a network whose dynamics are governed by constraint
equations analogous to the RSVP field equations.

(3)~Neural population coding in auditory cortex
implements something functionally equivalent to
\HCFone{} conformity encoding: correlated neural
populations are represented efficiently through
consensus encoding of shared dynamics rather than
independent encoding of individual firing rates.

Each of these three claims is at \LevelII{} (computational
metaphor) in its current formulation and would require
specific experimental designs to elevate to \LevelI{}.
The RSVP translation makes the experimental requirements
explicit by specifying what the RSVP-compatible versions
of these claims predict about measurable neural dynamics.

\medskip

Chapter~19 turns to the Nexus framework --- a new case
study that enters the monograph at the reconstruction
stage rather than the immersion stage. Unlike the Phoenix
corpus, which was the primary object of the audit, the
Nexus framework is introduced here as a parallel development
that instantiates the same RSVP principles the reconstruction
has been developing, but at the level of civilization-scale
computational infrastructure rather than individual
cognitive systems. Its inclusion is motivated by the
observation that the extensional--intensional distinction
central to the Phoenix critique applies with equal force
to AI computational memory as a whole.
